\documentclass[11pt]{article}
\usepackage[margin=1in]{geometry}
\usepackage{amsmath,amssymb}
\usepackage{booktabs}
\usepackage{graphicx}
\usepackage{hyperref}
\usepackage{xcolor}
\usepackage[numbers]{natbib}
\usepackage{enumitem}
\usepackage{caption}
\usepackage{subcaption}
\usepackage{multirow}
\setlist{nosep}

\graphicspath{{../results/}}

\newcommand{\Gr}{\mathrm{Gr}}

\title{Which geometric methods survive decoherence?\\
  A pre-registered comparison of 23 methods on\\
  simulated NV-diamond thermometry}
\author{Elliot Tower\textsuperscript{1}}
\date{}

\begin{document}
\maketitle

\noindent\textsuperscript{1}Independent Researcher\\
\noindent Correspondence: \href{mailto:elliot@elliottower.com}{elliot@elliottower.com}

\bigskip

% ================================================================
\begin{abstract}
Geometric methods for cross-sensor analysis---bracket norms, Grassmannian
geodesics, sheaf cohomology, discrete curvatures, persistent homology,
quantum information metrics---each detect specific structural properties
of multi-sensor data.  Which of these properties survive realistic
sensor noise, and which are artifacts of clean-data assumptions?  We
answer this for NV-diamond intracellular thermometry through a
pre-registered experiment comparing 23 geometric methods and 7 scalar
baselines across five decoherence axes (T2 degradation, surface noise,
device variation, temperature drift, ensemble size), with 20-seed
replication and bootstrap confidence intervals.  The central finding is
a method--axis sensitivity map: no single geometric method generalizes
across all decoherence types, and the axis that discriminates between
methods (device calibration variation) has the simplest physics
(additive offset).  We pre-registered that Grassmannian geodesic
distance would outperform bracket norms at detecting subspace drift;
it does not (median $\Delta\rho = -0.38$ on the composite axis,
$-0.74$ on the clean axis, both CIs excluding zero in the wrong
direction).  The refutation has a geometric explanation: geodesic
distance is invariant to mean shifts by construction, making it blind
to the dominant real-world noise source.  Of three confirmatory
hypotheses, one is confirmed (von Neumann entropy tracks decoherence
monotonically, $|\rho| = 0.90$) and two are refuted with mechanism.
The method--axis map, the device-variation discriminating axis, and the
subspace-method invariance result are the paper's contributions.
\end{abstract}

\noindent\textbf{Keywords:}
geometric methods, quantum sensors, NV-diamond, decoherence,
pre-registration, Grassmannian, bracket norm, persistent homology,
cross-sensor robustness

% ================================================================
\section{Introduction}
\label{sec:intro}

When multiple sensors measure the same quantity, they should agree.
Detecting disagreement requires a metric, and geometric methods offer a
rich toolkit: Grassmannian geodesic distance between PCA
subspaces~\citep{edelman1998geometry}, bracket norms measuring
inter-group variability in Lie-algebraic
terms~\citep{tower2026bracket}, sheaf cohomology
detecting device-level inconsistency~\citep{hansen2019toward},
discrete curvatures characterizing graph
structure~\citep{ollivier2009ricci,forman2003bochner},
persistent homology tracking topological features across
scales~\citep{carlsson2009topology}, and quantum information
metrics (QFI, Berry phase, von Neumann entropy) measuring sensitivity
and mixedness of quantum states~\citep{braunstein1994statistical}.

Each method encodes a different geometric property.  The question this
paper addresses is empirical: which of these properties are robust to
the noise sources that dominate real sensor deployments, and which
degrade or become uninformative as decoherence increases?

We answer this for NV-diamond intracellular thermometry, a quantum
sensing modality in which nitrogen-vacancy centers in diamond detect
temperature via the zero-field splitting shift ($dD/dT = -74$~kHz/K).
The simulation sweeps five physically motivated decoherence axes: T2
coherence time degradation (a composite axis affecting T2, linewidth,
and fluorescence contrast simultaneously), surface noise from diamond
defects, device-to-device calibration variation, temperature drift, and
ensemble size.  Each axis isolates a distinct noise mechanism, enabling
us to map which geometric methods respond to which noise type.

The experiment is pre-registered: all 23 methods, 7 baselines,
hypotheses, and statistical analysis plan were frozen under commit SHA
\texttt{8401d54} before any data was generated.  Two rounds of
adversarial review identified and corrected five structural problems
(insufficient replication, missing multiple-testing correction, composite
axis confounding, broken TDA implementations, unfair cross-method
comparison) before the freeze.  The full pre-registration, deviation
log, and raw results (30,000 method evaluations) are publicly
available.\footnote{\url{https://github.com/elliottower/geometric-sensor-robustness}}

\paragraph{Contributions.}
\begin{enumerate}
  \item A \textbf{method--axis sensitivity map} showing which geometric
    families respond to which decoherence types, revealing that no method
    generalizes across all five axes.
  \item A \textbf{pre-registered refutation} of the hypothesis that
    Grassmannian geodesic distance outperforms bracket norms, with a
    geometric explanation (mean-shift invariance).
  \item The identification of \textbf{device calibration variation} as
    the discriminating axis: 25 of 30 methods track the composite
    coherence axis at $|\rho| > 0.5$, but only 6 do so on device
    variation---and zero exceed $|\rho| > 0.8$.
\end{enumerate}

% ================================================================
\section{Simulation oracle}
\label{sec:oracle}

The NV-diamond simulation generates multi-device datasets with controlled
decoherence.  Each dataset contains $n = 200$ samples per class (healthy
tissue at 310.0~K, tumor tissue at 311.5~K) measured across $k = 3$
devices, yielding features derived from the ODMR (optically detected
magnetic resonance) spectrum: resonance frequency, linewidth, contrast,
T2 coherence time, and signal-to-noise ratio.

Five decoherence axes parameterize noise (Table~\ref{tab:axes}).  The
experimental design sweeps one axis at a time (10 levels each) while
holding the others at baseline (zero decoherence), with 20 independent
seeds per condition.

\begin{table}[t]
\centering
\caption{Decoherence axes.  \textbf{Composite axes} affect multiple
  features simultaneously; results on these axes must be interpreted
  with the coupling in mind.  sigma\_device is the cleanest axis
  (additive offset only) and serves as the primary discriminator
  between methods.}
\label{tab:axes}
\begin{tabular}{@{}llccl@{}}
\toprule
Axis & Parameter & Range & Composite? & Physics \\
\midrule
T2 degradation & $\alpha_{T2}$ & $[0, 1]$ & \textbf{Yes} & T2 + linewidth + contrast \\
Surface noise & $\sigma_\text{surface}$ & $[0, 0.5]$ & Yes & T2 noise + linewidth variance \\
Device variation & $\sigma_\text{device}$ & $[0, 0.1]$ & No & Additive temperature offset \\
Temperature drift & $\sigma_\text{temp}$ & $[0, 2.0]$~K & No & Per-sample noise \\
Ensemble size & $n_\text{centers}$ & $\{5 \ldots 50\}$ & No & Dimensionality change \\
\bottomrule
\end{tabular}
\end{table}

% ================================================================
\section{Methods}
\label{sec:methods}

All 23 geometric methods and 7 baselines take identical input
$(X, y, \text{device\_ids})$ and return a scalar score.  Methods are
grouped by geometric family (Table~\ref{tab:methods}).

\begin{table}[t]
\centering
\caption{Methods by family.  ``New'' indicates methods implemented for
  this study; others are ported from existing repositories.  Method~24
  (Chern number) was cut pre-registration due to a gauge-invariance
  error; method~15 is restricted to H0 (connected components) after the
  H1 heuristic was found to be incorrect.}
\label{tab:methods}
\small
\begin{tabular}{@{}llp{7cm}@{}}
\toprule
Family & Methods & Description \\
\midrule
Bracket (5) & 1--5 & Raw, corrected, transport-stable, phenotype-projected, localized bracket norms \\
Curvature (2) & 6--7 & Ollivier-Ricci (metric), Forman-Ricci (combinatorial) \\
Grassmannian (2) & 8--9 & Geodesic distance, loop holonomy \\
Sheaf (3) & 10--11, 23 & Cochran's Q via sheaf H$^1$, per-edge Q, persistent sheaf cohomology \\
Topology (2) & 15--16 & Persistent H0 (union-find), Wasserstein persistence distance \\
Quantum (4) & 17--18, 21--22 & QFI flatness, Berry phase proxy, von Neumann entropy stability, fidelity proxy \\
Alignment (2) & 13--14 & Linear CKA, Procrustes distance \\
Transport (2) & 12, 20 & TransportKit fusion, Wasserstein-1 distance \\
\midrule
\multicolumn{3}{@{}l}{\textit{Baselines (7):} t-test, Cohen's d, LASSO (3-fold CV AUC), RF, logistic, IRM penalty, DRO worst-device AUC} \\
\bottomrule
\end{tabular}
\end{table}

\subsection{Statistical framework}

All cross-method comparisons use Spearman rank correlation between a
method's score and the decoherence level, computed per seed and
summarized as the median $\rho$ across 20 seeds with bootstrap 95\% CI
(10,000 resamples).  This normalizes across methods' different score
scales by reducing each to a monotonic association strength on $[-1, +1]$.

Directional comparisons compute the per-seed paired difference
$\rho_A - \rho_B$ and report the bootstrap CI on the median difference.
If the CI excludes zero, the comparison is significant.

\paragraph{Confirmatory vs.\ exploratory.}
One correctness gate (H4d: sheaf H$^1$ = Cochran's Q, analytic
identity, no $\alpha$ needed) runs first.  The pre-registration froze
four confirmatory tests at $\alpha = 0.05/4 = 0.0125$: H3a on
$\alpha_{T2}$, H3a on $\sigma_\text{device}$ (two tests for the same
hypothesis on different axes), H6d, and H8a.  Post-execution, H8a was
voided as untestable under the Spearman framework (direct cross-scale
score comparison is incompatible with the normalization; see
Section~\ref{sec:deviation-h8a}), reducing the denominator from 4 to 3
and yielding a final $\alpha = 0.05/3 = 0.0167$.  This change is logged
in the deviation record.  All other hypotheses are exploratory.

% ================================================================
\section{Results}
\label{sec:results}

\subsection{Correctness gate}

H4d passes: the sheaf H$^1$ implementation reproduces Cochran's Q to
machine precision on synthetic data with known heterogeneity (5 devices,
3 with effect, 2 without; $Q = 168.1$, $p < 10^{-6}$) and returns
$Q \approx 0$ on homogeneous data ($Q = 0.28$, $p = 0.99$).

\subsection{Confirmatory hypotheses}

\paragraph{H3a: Grassmannian geodesic vs.\ bracket norm (REFUTED).}
On $\alpha_{T2}$ (composite axis), the median paired difference
$\rho_\text{geodesic} - \rho_\text{bracket}$ is $-0.38$ (95\% CI:
$[-0.48, -0.32]$).  On $\sigma_\text{device}$ (clean axis), the
difference is $-0.74$ ($[-0.94, -0.52]$).  Both CIs exclude zero in
the wrong direction: bracket norm strictly outperforms Grassmannian
geodesic.

The level-by-level data reveals the mechanism.
On $\sigma_\text{device}$, mean geodesic scores are flat between 1.94
and 1.95 across all decoherence levels, with no monotonic trend
($\rho = -0.15$, CI includes zero).
Device calibration variation adds a constant temperature offset to all
samples from a given device.  This shifts both class means equally,
preserving the principal angle between class subspaces.  The
Grassmannian geodesic is geometrically correct in reporting ``no
subspace change''---but this invariance makes it blind to the dominant
noise source in multi-device deployments.

Bracket norms, by contrast, respond to both subspace rotation and
scale changes in inter-group variability, achieving $\rho = +0.70$
on $\sigma_\text{device}$.

\paragraph{H6d: Von Neumann entropy monotonicity (CONFIRMED).}
Von Neumann entropy stability achieves $|\rho| = 0.90$ on
$\alpha_{T2}$ (CI: $[0.87, 0.93]$), exceeding the pre-registered
0.8 threshold at $\alpha = 0.0167$.

\paragraph{H8a: LASSO vs.\ geometric at low decoherence (VOIDED).}
\label{sec:deviation-h8a}
This hypothesis required comparing LASSO AUC directly against geometric
method scores at low decoherence ($\alpha_{T2} < 0.22$).  At these
levels, all classification baselines achieve AUC = 1.0 on all 20 seeds
(the 1.5~K tumor-healthy difference produces a $\sim$111~kHz frequency
shift that remains trivially classifiable).  Geometric methods produce
scores on incomparable scales (bracket norms $\sim 10^8$, curvatures
$\sim -20$, angles $\sim 2.0$).  The Spearman-$\rho$ normalization
adopted for all other comparisons is inapplicable to a within-regime
score comparison.  The underlying observation---classification is
trivially solved at low decoherence while geometric methods detect
structural changes below the classification threshold---is reported
descriptively but carries no confirmatory weight.

\subsection{The method--axis sensitivity map}

Table~\ref{tab:sensitivity} presents the core result: the median
Spearman $\rho$ for each method family's best representative on each
decoherence axis.  The ``Axes $> 0.8$'' column counts how many axes
the method tracks with $|\rho| > 0.8$.

\begin{table}[t]
\centering
\caption{Method--axis sensitivity map.  Each cell shows the median
  Spearman $\rho$(score, decoherence level) across 20 seeds for the
  best method in each family.  Bold indicates $|\rho| > 0.8$.
  $\sigma_\text{dev}$ is the discriminating axis: most methods are near
  zero.  $\dagger$Classification baselines return AUC = 1.0 (constant)
  on $\sigma_\text{surf}$ and $\sigma_\text{dev}$; Spearman $\rho$ is
  undefined.}
\label{tab:sensitivity}
\small
\begin{tabular}{@{}llcccccr@{}}
\toprule
Family & Best method & $\alpha_{T2}$ & $\sigma_\text{surf}$ & $\sigma_\text{dev}$ & $\sigma_\text{temp}$ & $n_\text{cent}$ & $>$0.8 \\
\midrule
Bracket & bracket\_norm\_raw & \textbf{+1.00} & \textbf{+0.90} & +0.70 & \textbf{+1.00} & \textbf{+0.96} & 4/5 \\
Bracket & transport\_stable & \textbf{+1.00} & \textbf{+0.94} & $-$0.10 & \textbf{+1.00} & \textbf{+0.98} & 4/5 \\
Topology & persistent\_h0 & \textbf{+1.00} & \textbf{+1.00} & +0.55 & \textbf{+0.98} & \textbf{+1.00} & 4/5 \\
Alignment & CKA & \textbf{$-$1.00} & \textbf{$-$0.98} & +0.12 & \textbf{$-$1.00} & \textbf{+1.00} & 4/5 \\
Quantum & fidelity\_proxy & \textbf{$-$0.96} & \textbf{$-$0.99} & $-$0.03 & \textbf{+1.00} & \textbf{$-$1.00} & 4/5 \\
Alignment & Procrustes & \textbf{+1.00} & \textbf{+1.00} & $-$0.02 & \textbf{+0.99} & \textbf{$-$1.00} & 4/5 \\
Sheaf & persistent\_sheaf & \textbf{$-$1.00} & \textbf{$-$1.00} & $-$0.16 & \textbf{$-$0.99} & $-$0.03 & 3/5 \\
\midrule
Grassmannian & geodesic & +0.62 & $-$0.52 & $-$0.15 & $-$0.16 & \textbf{+0.87} & 1/5 \\
Curvature & Ollivier-Ricci & $-$0.59 & +0.75 & +0.07 & $-$0.46 & \textbf{+0.93} & 1/5 \\
Sheaf & sheaf\_h1 & $-$0.29 & $-$0.21 & $-$0.07 & +0.20 & $-$0.67 & 0/5 \\
\midrule
Classification$^\dagger$ & LASSO & \textbf{$-$0.94} & --- & --- & \textbf{$-$0.99} & --- & 2/5 \\
Effect size & Cohen's d & \textbf{$-$1.00} & $-$0.25 & $-$0.72 & \textbf{$-$1.00} & $-$0.11 & 2/5 \\
\bottomrule
\end{tabular}
\end{table}

\subsection{Sigma\_device: the discriminating axis}

Of 30 methods tested, only 6 achieve $|\rho| > 0.5$ on
$\sigma_\text{device}$: localized bracket ($-0.74$), Cohen's d
($-0.72$), bracket norm raw ($+0.70$), bracket norm corrected
($+0.70$), persistent H0 ($+0.55$), and IRM penalty ($-0.52$).
None exceed $|\rho| > 0.8$.
By contrast, 25 of 30 methods achieve $|\rho| > 0.5$ on
$\alpha_{T2}$, and 20 exceed $|\rho| > 0.8$.

The physics explains the gap.  Device calibration variation adds a
constant temperature offset per device---a mean shift in the feature
space that preserves within-device covariance structure.  Methods
operating on subspace angles (Grassmannian), graph curvature
(Ollivier-Ricci, Forman-Ricci), spectral properties, or
direction-based proxies are invariant to such shifts by construction.
Only methods sensitive to between-group location differences
(bracket norms, effect sizes) or within-group dispersion changes
(persistent H0) detect the perturbation.

This result has practical implications for clinical multi-sensor
deployments, where device-to-device calibration differences are often
the primary noise source.  Geometric methods that are invariant to
additive offsets provide no advantage over a t-test in this regime.

\subsection{Exploratory findings}

\paragraph{Direction-based methods share mean-shift invariance.}
Berry phase, Grassmannian geodesic, and Grassmannian holonomy all
measure angles or direction-dependent properties in feature space.
All three are blind to $\sigma_\text{device}$ ($|\rho| < 0.15$) for
the same reason: additive offsets do not change directions.  The
Berry-phase proxy in particular achieves $\rho = +1.00$ on
$\alpha_{T2}$ and $+0.95$ on $\sigma_\text{surface}$---both axes
that distort covariance structure---while remaining at $\rho = +0.006$
on $\sigma_\text{device}$.  This selectivity reflects the shared
geometric property (direction invariance to mean shifts), not a
quantum-specific detection mechanism.

\paragraph{Curvature phase transition.}
Ollivier-Ricci and Forman-Ricci curvature exhibit a sharp phase
transition on $\alpha_{T2}$: Ollivier-Ricci collapses from $0.998$ to
$0.200$ between $\alpha_{T2} = 0.33$ and $0.44$, reaching $\approx 0$
by $0.56$; Forman-Ricci jumps from $-18$ (negative curvature,
indicating structured clusters) to $+0.5$ at the same boundary.  This
transition marks the decoherence level at which the k-NN graph loses
cluster structure, defining a ``graph collapse threshold'' beyond which
discrete curvature methods are uninformative.

\paragraph{Persistent H0: the silent generalist.}
Persistent H0 (connected components via union-find) achieves 4/5 axes
at $|\rho| > 0.8$ with monotonically increasing total persistence on
$\alpha_{T2}$ (from $2.2 \times 10^7$ to $4.2 \times 10^8$).  The
decision to restrict to H0 and cut the incorrect H1 heuristic is
validated: the NV-diamond point cloud has no decoherence-varying loop
structure, so H1 would have added noise.

\paragraph{Sheaf cohomology: a split result.}
Sheaf H$^1$ (Cochran's Q) and per-edge sheaf Q achieve $|\rho| < 0.3$
on all axes---an honest null.  Persistent sheaf cohomology, by contrast, achieves
$|\rho| > 0.8$ on 3 of 5 axes ($\alpha_{T2}$: $-1.00$;
$\sigma_\text{surface}$: $-1.00$; $\sigma_\text{temp}$: $-0.99$),
while remaining blind to $\sigma_\text{device}$ ($-0.16$) and
$n_\text{centers}$ ($-0.03$).  Persistent sheaf cohomology combines
sheaf consistency checks with a persistence filtration, capturing
topological changes in the device-consistency structure across scales
that the pointwise Cochran Q misses.

A device-count expansion confirms that the sheaf H$^1$ null on
$\sigma_\text{device}$ is structural, not a power limitation.
At 5 devices (20 seeds), sheaf H$^1$ achieves $\rho = +0.04$
(CI: $[-0.16, +0.21]$)---statistically indistinguishable from the
3-device result ($\rho = -0.07$, CI: $[-0.22, +0.35]$).  By
contrast, bracket norm raw strengthens from $+0.70$ at 3 devices
to $+0.73$ at 5 devices (CI: $[+0.69, +0.76]$), indicating that
power increases with device count for methods sensitive to mean
shifts.  The Cochran Q statistic underlying sheaf H$^1$ measures
heterogeneity in per-device effect sizes; when all devices share
the same physics (single-modality NV-diamond), a common-mode
calibration offset produces homogeneous effects by construction,
yielding $Q \approx 0$ regardless of device count.

\paragraph{Spectral gap: an $\alpha_{T2}$ specialist.}
Spectral gap stability achieves $\rho = +0.96$ on $\alpha_{T2}$ and
$|\rho| < 0.12$ on all other axes.  The composite coherence axis
collapses the spectral gap of the sample covariance matrix, while
the other four axes (additive offsets, per-sample noise, dimensionality)
preserve the eigenvalue gap.  This makes spectral gap stability a
narrow but precise detector of coherence loss.

\paragraph{Localized bracket: a sigma\_device specialist.}
Localized bracket achieves $\rho = -0.74$ on $\sigma_\text{device}$
and $|\rho| < 0.05$ on all other axes.  It measures within-neighborhood
geometric structure, which collapses as device offsets increase the
intra-device spread.  This is the only method with a single-axis
specialization profile.

\paragraph{Classification baselines saturate at low decoherence.}
All classification baselines (LASSO, logistic, RF, DRO) return AUC =
1.0 on every level of $\sigma_\text{surface}$ and
$\sigma_\text{device}$ across all 20 seeds.  The 1.5~K tumor-healthy
temperature difference ($\sim$111~kHz frequency shift) remains
trivially classifiable even at maximum surface noise or device
variation.  On $\alpha_{T2}$, AUC begins degrading above
$\alpha_{T2} \approx 0.44$, with LASSO dropping from 1.0 to 0.75 at
$\alpha_{T2} = 1.0$.  Geometric methods detect structural changes at
decoherence levels where classification has not yet degraded---they
measure different properties than classification accuracy.

% ================================================================
\section{Discussion}
\label{sec:discussion}

\subsection{No universal geometric method}

The central finding is negative: no geometric method achieves
$|\rho| > 0.8$ on all five decoherence axes.  Six methods cover 4 of
5 axes (bracket norm raw, transport-stable bracket, persistent H0,
CKA, Procrustes, fidelity proxy), and persistent sheaf cohomology
covers 3 of 5.  The missing axis for all of them is
$\sigma_\text{device}$ (or, for the transport-stable bracket and sheaf
methods, $n_\text{centers}$ as well).  Methods designed to capture subspace geometry
(Grassmannian), graph topology (curvatures), or direction-dependent
properties (Berry phase proxy) achieve strong performance on their
responsive axes but are blind to additive noise.

This pattern is not a failure of the methods.  Each method correctly
computes the geometric property it was designed to measure.  The
finding is that the geometric property relevant to decoherence depends
on the noise mechanism: coherence degradation rotates and scales
subspaces ($\alpha_{T2}$), surface defects perturb graph structure
($\sigma_\text{surface}$), but device variation merely shifts
locations ($\sigma_\text{device}$).  Matching the method to the noise
geometry is the operational challenge.

\subsection{The value of the companion axis}

The $\sigma_\text{device}$ companion axis was added during
pre-registration review to disentangle composite-axis effects from
genuine subspace drift.  Without it, the paper would have reported that
``most methods work well'' (25/30 achieve $|\rho| > 0.5$ on
$\alpha_{T2}$) and missed the central finding.  The companion axis design---testing the same
hypothesis on a clean, non-composite axis alongside the composite
one---is a transferable experimental principle for any study where the
primary axis affects multiple features simultaneously.

\subsection{Pre-registration discipline}

The H3a refutation demonstrates the value of pre-registration.  Post
hoc, one might rationalize either outcome: ``Grassmannian methods should
detect subspace drift because they directly measure subspace angles''
(our pre-registered prediction) or ``Grassmannian methods should miss
additive noise because they are location-invariant'' (the actual
explanation).  Both are geometrically coherent.  Without a frozen
prediction, the second explanation could be presented as if it were
the original hypothesis.

Similarly, the H8a voiding shows that pre-registration can expose
incoherence in the statistical framework.  The hypothesis was written
before the Spearman-$\rho$ normalization was adopted, and it survived
the review rounds because it seemed intuitively correct.  Only when the
results arrived did the scale-incompatibility become undeniable.

\subsection{Limitations}

The NV-diamond simulation is a controlled testbed, not a substitute for
real sensor data.  Feature dimensions are low ($\leq 5$), sample sizes
are moderate ($n = 200$ per class), and the noise model is parametric.
Real NV-diamond measurements involve nonlinear transductions,
non-Gaussian noise, and correlated decoherence channels.  The specific
$\rho$ values reported here depend on these simulation choices.

The sheaf H$^1$ null on $\sigma_\text{device}$ is structural, not a
power limitation: a device-count expansion from 3 to 5 devices (20
seeds each) shows no activation ($\rho$ shifts from $-0.07$ to
$+0.04$, both CIs including zero), while bracket norms strengthen
over the same range.  Single-modality device variation produces
homogeneous per-device effects, which Cochran's Q cannot detect by
construction.  Persistent sheaf cohomology, which wraps Q in a
persistence filtration, achieves strong performance on 3 of 5
axes---the split result reflects the difference between pointwise
heterogeneity testing and multi-scale topological structure.  Testing
sheaf H$^1$ on cross-modality data, where different sensor types
produce qualitatively different effect sizes, remains an open
question.

The Berry-phase and QFI implementations are classical proxies applied to
feature-space statistics, not direct quantum-state measurements.  Claims
about ``quantum-channel selectivity'' would require access to the actual
quantum state, not to processed measurement features.

% ================================================================
\section{Conclusion}
\label{sec:conclusion}

This study provides a pre-registered map of which geometric methods
respond to which decoherence mechanisms in simulated NV-diamond
thermometry.  The headline result is a refutation: Grassmannian geodesic
distance, a principled subspace metric, is blind to the noise source
most relevant to clinical multi-sensor deployments (device calibration
variation) because of a geometric invariance property.  Six methods (bracket norms, persistent H0, CKA, Procrustes,
fidelity proxy, transport-stable bracket) achieve the best
generalization at 4 of 5 axes, but all fall short on device variation.  The
method--axis sensitivity map produced here can guide method selection for
specific deployment scenarios: use bracket norms for broad coverage,
direction-based methods for coherence degradation detection, and simple
effect-size baselines when device variation dominates.

\paragraph{Reproducibility.}
All code, pre-registration, deviation log, and raw results (30,000
evaluations) are available at
\url{https://github.com/elliottower/geometric-sensor-robustness}.
The simulation is fully deterministic: given the frozen commit SHA
(\texttt{8401d54}) and the 20-seed pool (seeds 1000--1019), all results
can be regenerated.

\subsection*{Data availability}

Raw JSON results (all 30,000 evaluations) and analysis output
(bootstrap CIs, confirmatory test outcomes) are committed alongside
the source code.

\subsection*{Software availability}

Source code requires Python~3.11+, NumPy, SciPy, scikit-learn, and
tqdm (no GPU or proprietary software).  MIT license.

\subsection*{Author contributions}

E.T.\ conceived the study, designed and implemented all experiments and
methods, performed the statistical analysis, and wrote the manuscript.

\subsection*{Competing interests}

No competing interests were disclosed.

\subsection*{Grant information}

No grants were involved in supporting this work.

\bibliographystyle{unsrtnat}
\bibliography{refs}

\end{document}
